\documentclass{article}
\usepackage{graphicx} % Required for inserting images
\usepackage{amsmath}
\usepackage{amsfonts}                                                                   
\usepackage{amsthm}
\usepackage{amssymb}
\usepackage[utf8]{inputenc}
\usepackage[left=2cm, top=-5cm, right=2cm, bottom=2cm, includeheadfoot, headheight=130pt, a4paper]{geometry}
\usepackage[polish]{babel}
\usepackage[utf8]{inputenc}
\usepackage{polski}
\usepackage[T1]{fontenc}
\usepackage{longtable}
\usepackage{caption}
\usepackage{subcaption}
\usepackage{booktabs}
\usepackage{array}
\usepackage{float}
\usepackage{hyperref}
\usepackage{float}
\usepackage{longtable}

\usepackage[polish]{babel}
\usepackage[utf8]{inputenc}
\usepackage{polski}
\usepackage[T1]{fontenc}

\title{Algorytmy Metaheurystyczne\\Lista 3 - TSP A Selection of Genetic Algorithms}
\author{Rafał Włodarczyk}
\date{02.06.2026}

\begin{document}

\maketitle

\tableofcontents

\section{Model}

Model TSP zakłada graf pełny $G=(V,E)$ z wagami.

\section{Źródło dla danych}

Dane do programu są wzięte ze strony internetowej:
\begin{center}
\texttt{https://www.math.uwaterloo.ca/tsp/world/countries.html}
\end{center}
oraz są dostępne do pobrania w formacie \texttt{.tsp} za pomocą skryptu \texttt{get\_data.sh}.

\section{Idea Algorytmu Genetycznego}

Algorytm Genetyczny zgodnie z opisem na liście składa się z pięciu kroków
\begin{enumerate}
    \item Populacja Początkowa. Wykorzystałem permutację powstałą za pomocą MST.
    \item Ewaluacja i Selekcja. Selekcja odbywa się za pomocą przyspieszonej metody turniejowej.
    Wybierane są dwa dowolne elementy populacji (część losowa), a następnie brany jest większy z nich (część turniejowa).
    \item Krzyżowanie. Wykorzystane zostaną dwa algorytmy \texttt{PMX Crossover} oraz \texttt{OX Crossover}.
    \item Mutacja. Odbywa się za pomocą wybrania dwóch indeksów $x,y$, a następnie odwrócenia podpermutacji - 
    wykonanie inwersji.
    \item Warunek Końca. Ustalony na sztywną liczbę generacji. Zwracamy najlepszy z wyników
\end{enumerate}

Ponadto rozwinięty zostanie o dwa dodatkowe elementy

\begin{enumerate}
    \item Memetyka - Przyspieszenie \textit{dojrzewania} nowych elementów populacji. Zastosowany jest skrócony algorytm Tabu Search,
    z ustaloną niską maksymalną liczbą iteracji.
    \item Wyspowość - Podzielenie populacji na mniejsze niezależne podpopulacje, które po ustalonej liczbie generacji wymieniają się informacjami
    za pomocą specjalnej funkcji.
\end{enumerate}

\section{Zadanie 1 - Metody krzyżowania}

Wybrałem dwie metody krzyżowania:
\begin{enumerate}
    \item PMX (Partially Mapped Crossover) - Goldberg / Lingle, 1985. 
    \item OX (Order Crossover) - Davis, 1985
\end{enumerate}

\subsection{Metoda PMX}

Metoda PMX (Partially Mapped Crossover) to algorytm krzyżowania permutacji 
stworzony z myślą o TSP. Bierzemy dwie permutacje $p_1,p_2$, będziemy tworzyć dziecko $c$. Bierzemy dwa dowolne punkty $0\leq x\leq y \leq n$, następnie dzielimy zbiór na trzy części:
lewy (indeksy $0\dots x - 1$), środkowy (indeksy $x \dots y$), oraz prawy (indeksy $y+1 \dots n$). 
\begin{enumerate}
    \item Segment środkowy $(x,y)$ zostaje powielony z $p_1$ do $c$
    \item Znajdujemy konflikty - przenosimy elementy z permutacji $p_2$ do $c$, o ile już w $c$ ich nie ma.
    Musimy znaleźć mapowania elementów $k_1,\dots, k_l$ do $c$ spoza segmentu środkowego. Niech $k=k_i$ (dowolny brakujący element).
    Szukamy indeksu $i : p_2(i) = k$, znajdujemy wartość $m=p_1(i)$, a następnie szukamy dla jakiego $j$ $p_2(j)=m$.
    Gdy $j \notin (x,y)$ możemy umieścić element $c(j)=k$.
    Osiągamy dzięki temu zestaw mapowania, który pozwala nam umieścić $k$ w wolne miejsce $c$.
\end{enumerate}

\subsection{Metoda OX}

Metoda OX (Order Crossover) to algorytm krzyżowania permutacji dla TSP. Bierzemy dwie permutacje $p_1,p_2$, będziemy tworzyć dziecko $c$. Bierzemy dwa dowolne punkty $0\leq x\leq y \leq n$, następnie dzielimy zbiór na trzy części:
lewy (indeksy $0\dots x - 1$), środkowy (indeksy $x \dots y$), oraz prawy (indeksy $y+1 \dots n$). 
\begin{enumerate}
    \item Segment środkowy $(x,y)$ zostaje powielony z $p_1$ do $c$
    \item Dokonujemy cyklicznego wypełnienia zaczynając od indeksu $i=y+1$, sprawdzamy czy element $p_2(i)$ jest już w $c$,
    jeśli jest to dodajemy indeks i przechodzimy dalej. Jeśli nie ma to wpisujemy $c(i)=p_2(i)$ i przesuwamy $i$ na następne wolne miejsce.
\end{enumerate}

\subsection{Wyniki porównania metod PXM i OX}

Dla zadanych danych OX jest lepszy od PMX. Zostanie zatem wykorzystany w zmodyfikowanych algorytmach.

\section{Zadanie 2 - Algorytm Wyspowy}

Algorytm wyspowy otrzymuje dodając dwa dodatkowe parametry
\begin{enumerate}
    \item \texttt{icount} - Ilość wysp, na które zostanie podzielony zbiór populacji.
    \item \texttt{interval} - Ilość iteracji po których zachodzi wymiana informacji pomiędzy wyspami.
    Niech $I$ to będzie nasz interwał. Wtedy do wymiany informacji dojdzie w iteracjach $k\cdot I, k\in\mathbb{N}_{\geq 0}$
\end{enumerate}

\section{Zadanie 3 - Wykorzystanie OpenMP w celu przyspieszenia obliczeń}

Wykorzystam bibliotekę OpenMP \texttt{<omp.h>} która pozwala na zrównoleglenie obliczeń - wykorzystuję 
rozdzielając część pętli na wiele rdzeni. Dzięki ustaleniu generatora liczb losowych per wątek mogę
uzyskiwać różne wyniki symulując konsekwentne wykonywanie kolejnych iteracji algorytmów.  

\section{Zadanie 4 - Algorytm Memetyczny}

Moja realizacja algorytmu memetycznego opiera się o dodanie Tabu Search po mutacji dzieci poprzedniej populacji.
Dzięki temu zdecydowanie szybciej (w obrębie jednej generacji) będziemy optymalizować
do kolejnego minimum lokalnego po mutacji - czyli po potencjalnym wyskoczeniu z minimum lokalnego.

\section{Porównanie Algorytmów z Listy 3}

\input{parameters.tex}

\section{Wykresy najlepszych rozwiązań}

\input{plots.tex}

\section{Finalne zestawienie wszystkich algorytmów}

\end{document}